\documentclass[12pt]{article}
\usepackage{graphicx}
\usepackage{geometry}
\usepackage{listings}
\usepackage{xcolor}

\usepackage[T1]{fontenc}

% --- Helpers for line breaks, floats, captions, and outlines ---
\usepackage{url}
\usepackage[hidelinks]{hyperref}
\usepackage{float}      % provides [H]
\usepackage{placeins}   % provides \FloatBarrier
\usepackage{caption}
\usepackage{bookmark}
\captionsetup{justification=raggedright,singlelinecheck=false}
\sloppy
\setlength{\emergencystretch}{3em}

\lstset{
  basicstyle=\ttfamily\small,
  breaklines=true,
  breakatwhitespace=true,
  columns=fullflexible
}

\geometry{margin=1in}

% ---- Macro for bottom-only cropping ----
% Usage: \cropbottom[<bottom>]{<file>}
% Example: \cropbottom[150pt]{step1.png}
\newcommand{\cropbottom}[2][150pt]{%
  \includegraphics[width=0.9\textwidth, trim={0 #1 0 0}, clip]{#2}%
}

\begin{document}

\begin{center}
    {\Large \textbf{COMP 324 Network Security}}\\[4pt]
    {\large Module 3 Lab 2: File Operations and Permissions}\\[10pt]
    \textbf{Lab Report}\\[6pt]
    \textbf{Student: Erik Gonzalez}\\[4pt]
    \textbf{Professor: Samuel Addington}\\[4pt]
    \textbf{Date: \today}
\end{center}



\section*{Exact Commands Run}
\begin{lstlisting}
# Step 1: Creating Files and Directories
mkdir -p /home/student1/project
touch /home/student1/project/test.txt
echo "This is a test file for access control lab" > /home/student1/project/test.txt
ls -la /home/student1/project/

# Step 2: Understanding File Permissions
ls -la /home/student1/project/test.txt

# Step 3: Changing File Ownership
sudo chown root:root /home/student1/project/test.txt
ls -la /home/student1/project/test.txt

# Step 4: Testing Access as Student
su - student1 -c "echo 'test' >> /home/student1/project/test.txt"
\end{lstlisting}


\section*{Step 1: Creating Files and Directories}
\begin{figure}[H] % force placement here
    \centering
    \cropbottom[1100pt]{step1.png}
    \end{figure}

\FloatBarrier % prevent floats crossing into the next section

\section*{Step 2: Understanding File Permissions}
\begin{figure}[H] % force placement here
    \centering
    \cropbottom[1250pt]{step2.png}
    
\end{figure}

\noindent \textbf{Recorded output:}
\begin{lstlisting}
[student1@threepio ~]$ ls -la /home/student1/project/test.txt
-rw-r--r-- 1 student1 student1 43 Sep 28 21:47 /home/student1/project/test.txt
\end{lstlisting}

\noindent \textbf{\\Field Breakdown}

\begin{itemize}
\item \textbf{File permissions:} 
\mbox{\texttt{%
\textcolor{black}{-}%
\textcolor{blue}{rw-}%
\textcolor{green}{r--}%
\textcolor{red}{r--}%
}} 
\mbox{(\textcolor{black}{regular file}, 
\textcolor{blue}{owner: read/write}, 
\textcolor{green}{group: read only}, 
\textcolor{red}{others: read only})}


    \item \textbf{Hard link count:} \texttt{1} → only one file name points to this file.

    \item \textbf{User owner:} \texttt{student1} → the user that owns the file.

    \item \textbf{Group owner:} \texttt{student1} → the group with access permissions.

    \item \textbf{File size (bytes):} \texttt{43} → the file contains 43 bytes of data.

    \item \textbf{Last modified date/time:} \texttt{Sep 28 21:47} → when the file was last changed.

    \item \textbf{Filename:} \url{/home/student1/project/test.txt} → the file’s full path.
\end{itemize}


\FloatBarrier % keep Step 2 figure/output with Step 2

\section*{Step 3: Changing File Ownership}
\begin{figure}[H] % force placement here
    \centering
    \cropbottom[400pt]{step3.png}
    \end{figure}


\noindent \textbf{\\Updated ownership:}\\
\texttt{-rw-r--r-- 1 root root 43 Sep 28 21:47 /home/student1/project/test.txt}


\FloatBarrier

\section*{Step 4: Testing Access as Student}
\begin{figure}[H] % force placement here
    \centering
    \cropbottom[400pt]{step4.png}
    
\end{figure}

\noindent \textbf{Recorded error message:}
\begin{lstlisting}
-bash: line 1: /home/student1/project/test.txt: Permission denied
\end{lstlisting}

\noindent \textbf{\\Error Explanation:}\\
This occurred because the file is now owned by root, and student1 only has read access, not write access.



\section*{Observations:}
\begin{itemize}
    \item File ownership decides which user can read, write, or modify a file.
    \item Changing ownership with \texttt{chown} can remove the original user’s ability to edit their own file.
    \item Permissions \texttt{rw-r--r--} allow only the owner to write, while group and others have read-only access.
    \item Attempting to append text as a non-owner is prohibited by Linux security.
 \end{itemize}
\section*{Conclusion}
File ownership and permissions are very important in Linux.  
This lab showed how creating files, changing ownership, and testing access shows how serious Linux is about it's permisions.  
Only the owner or root can modify files with no resctriction, while others are restricted.  \end{document}

